% GENERATED FILE. Edit canonical lessons, then run npm run generate:books.
% canonical-source-hash: fnv1a64:994167f0465e7203
% canonical-lessons: ES-C462-reducir, ES-C462-decidirse, ES-C462-jefe, ES-R462-repaso-jefe, ES-C462-sintesis-jefe

\chapter{La reunión}
\label{ch:es-jefe}
\hlchaptermodality{462}

\begin{chapteropening}
\textbf{By the end of this chapter:} \emph{I can follow a message from a manager about cutting spending, and reply that I am still choosing between two options.}

Read the message or notice this chapter's words exist for, then say the same thing back.
\end{chapteropening}

\section[reducir]{reducir — to bring it down}

\label{lesson:ES-C462-reducir}



\begin{quote}
Say \textbf{la ducha} and \textbf{el coste}. A shower is about to turn out to be a relative of the verb coming next.
\end{quote}

\subsection*{You'll want to know: reducir}
\textbf{reducir} — to reduce, to cut back.

\emph{Reducir costes} — to cut costs, and \emph{reducir costos} wherever you met \emph{el costo}. \emph{Reducir la velocidad} — to slow down, the wording on a road sign. \emph{Han reducido la plantilla} — they've cut the staff.

\textbf{La reducción} is the noun, and \emph{reducido} on a price tag means marked down.

\begin{grammarlens}[title={Grammar Lens: the -zc- verbs}]
The \textbf{c} is already in \emph{reducir}. What this group grows is a \textbf{z} in front of it:

\noindent
\begin{tabularx}{\linewidth}{@{}>{\raggedright\arraybackslash}X>{\raggedright\arraybackslash}X@{}}
\toprule
\textbf{verb} & \textbf{I …} \\
\midrule
reducir & redu\textbf{zc}o \\
conducir & condu\textbf{zc}o \\
conocer & cono\textbf{zc}o \\
\bottomrule
\end{tabularx}

It is not only the \emph{yo} form. The whole subjunctive carries the \textbf{z} too — \emph{reduzca, reduzcas, reduzcamos} — and that is where you will actually read it, because the polite command is the subjunctive: \textbf{Reduzca la velocidad}, the road sign from a moment ago. The rest of the present stays plain: \emph{reduces}, \emph{reduce}, \emph{reducimos}.

Watch the past, though, because the table above hides a split. \emph{Reducir} and \emph{conducir} go \emph{reduje} and \emph{conduje} — not \emph{reducí}. \emph{Conocer} shares the \textbf{-zc-} and not that past: it goes \emph{conocí}. Two groups, one spelling.

The \textbf{z} has two different stories as well. In \emph{conocer} it is inherited: Latin \emph{cognoscere} had the cluster and the \emph{yo} form kept it. Latin \emph{reducere} never had it, so \emph{reduzco} got its \textbf{z} by copying the verbs that did — a pattern that spread rather than a sound that survived. The family it spread through is large, several hundred verbs once the \emph{-ecer} verbs are counted in.
\end{grammarlens}

\begin{cousinweb}
Latin \textbf{reducere}: \textbf{re-}, back, plus \textbf{ducere}, \textbf{to lead}. To reduce something is to \textbf{lead it back} — to a smaller size, to where it was.

\emph{Ducere} is everywhere once you see it. \textbf{La ducha} is the same root, and it travelled: Latin \emph{ductia} gave Italian \emph{doccia}, which gave French \emph{douche}, which gave Spanish \emph{ducha} — a pipe that \textbf{leads} water. English has \textbf{conduct}, \textbf{produce}, \textbf{educate} (to lead out), \textbf{duke} — a leader — and \textbf{aqueduct}, which leads water exactly as a shower does.

Spanish keeps \emph{conducir}, to drive, and \emph{el conductor}. A driver leads the car; a shower leads the water; reducing leads the number back down.
\end{cousinweb}

\subsection*{Your turn}
Say these aloud:

\begin{itemize}
\raggedright
  \item we have to cut costs, in the variety you use
  \item I reduce, using the first-person form
  \item what a shower and a duke have in common
\end{itemize}

\begin{tcolorbox}[breakable,colback=teal!4,colframe=teal!35!black,title={Before you move on}]
To cut back? (\textbf{Reducir}.) I reduce? (\textbf{Reduzco}.) What does \emph{ducere} mean? (\textbf{To lead}.)
\end{tcolorbox}
\section[decidirse]{decidirse — to make your mind up}

\label{lesson:ES-C462-decidirse}



\begin{quote}
Say \textbf{decidir} and \textbf{acercar}. You met \emph{decidirse} in passing when you learned \emph{decidir}; now take it properly, with what \textbf{-se} did to the second word in front of you.
\end{quote}

\subsection*{You'll want to know: decidirse}
\textbf{decidirse} — to make up your mind.

\emph{No me decido} — I can't decide. \emph{Decídete} — make your mind up. \emph{Se decidió por el azul} — she went for the blue one.

The \textbf{-se} version takes \textbf{por} for the thing chosen: \emph{me decidí por el piso pequeño}.

\begin{grammarlens}[title={Grammar Lens: what the -se adds here}]
You have met the \emph{-se} switch three times, and this is a fourth use of it — not a new meaning, but a new \emph{shade}:

\noindent
\begin{tabularx}{\linewidth}{@{}>{\raggedright\arraybackslash}X>{\raggedright\arraybackslash}X@{}}
\toprule
\textbf{without -se} & \textbf{with -se} \\
\midrule
\emph{acercar} — bring it closer & \emph{acercarse} — come closer \\
\emph{decidir} — settle a matter & \emph{decidirse} — settle \textbf{yourself} \\
\bottomrule
\end{tabularx}

\emph{He decidido vender el coche} states the decision. \emph{Me he decidido a vender el coche} says you struggled and got there. The reflexive adds the \textbf{effort}, the sense of having had to bring yourself to it.

That is why \emph{no me decido} is the natural way to say you are dithering, and \emph{no decido} would sound like a job description.
\end{grammarlens}

\begin{cousinweb}
Latin \textbf{decidere}: \textbf{de-}, off, plus \textbf{caedere}, \textbf{to cut}. To decide is to \textbf{cut off} — the other options, the deliberating, the retreat.

That picture explains the English family exactly. A \textbf{decision} cuts things off. Something \textbf{decisive} cuts cleanly. And \emph{caedere} is also in \textbf{concise} (cut short), \textbf{precise} (cut in front, trimmed to the edge), and — through the same cutting — \textbf{homicide} and \textbf{suicide}.

Spanish has \emph{decidir}, \emph{la decisión}, \emph{conciso} and \emph{preciso} off the same root. When you make your mind up, a Latin speaker heard you reaching for a blade.
\end{cousinweb}

\subsection*{Your turn}
Say these aloud:

\begin{itemize}
\raggedright
  \item I can't make my mind up
  \item she went for the blue one
  \item which preposition follows \emph{decidirse}
\end{itemize}

\begin{tcolorbox}[breakable,colback=teal!4,colframe=teal!35!black,title={Before you move on}]
To make your mind up? (\textbf{Decidirse}.) What does the \emph{-se} add? (\textbf{The effort of getting there}.) What does \emph{caedere} mean? (\textbf{To cut}.)
\end{tcolorbox}
\section[el jefe]{el jefe — the one at the head}

\label{lesson:ES-C462-jefe}



\begin{quote}
Say \textbf{la cabeza} and \textbf{el trabajo}. Put the first at the top of the second and you have named somebody.
\end{quote}

\subsection*{You'll want to know: el jefe}
\textbf{el jefe} — the boss, the head of something.

\emph{Hablar con el jefe} — to speak to the boss. \emph{La jefa de personal} — the head of HR. \emph{El jefe de estación} — the stationmaster.

The feminine is \textbf{la jefa}, formed normally, and both are ordinary words at work — not slang and not deferential.

\begin{grammarlens}[title={Grammar Lens: the j that used to be something else}]
\textbf{Jefe} came from French \emph{chef}, and the \textbf{ch} became a \textbf{j}. Spanish did that to a whole shelf of borrowings:

\noindent
\begin{tabularx}{\linewidth}{@{}>{\raggedright\arraybackslash}X>{\raggedright\arraybackslash}X@{}}
\toprule
\textbf{French} & \textbf{Spanish} \\
\midrule
\textbf{ch}ef & el \textbf{j}efe \\
\textbf{j}ardin & el \textbf{j}ardín \\
\textbf{j}amais & — \\
\bottomrule
\end{tabularx}

The \textbf{j} is the \emph{jota}, the sound at the back of the throat, and it is a latecomer: medieval Spanish did not have it. Words spelled with \textbf{j} or with \textbf{g} before \emph{e} and \emph{i} all shifted to it at roughly the same time, a few hundred years ago, which is why \emph{jefe}, \emph{mujer} and \emph{gente} now share a sound that none of them started with.
\end{grammarlens}

\begin{cousinweb}
French \textbf{chef} is Latin \textbf{caput}, \textbf{head} — and \emph{caput} is the same word inside \textbf{la cabeza}.

That needs one correction to something you were told earlier. When you learned \emph{la cabeza}, French was described as having dropped \emph{caput} altogether for \emph{testa}, a pot. That is true of the \textbf{body part} and only of it: \emph{tête} replaced \emph{caput} for the thing on your shoulders, while \emph{caput} itself stayed on in French as \emph{chef}, a head in the sense of a chief. It is that surviving branch Spanish later borrowed.

So \emph{el jefe} and \emph{la cabeza} are one Latin noun arriving twice: once worn down through Spanish speech into the body part, once borrowed back from French as the person in charge. A boss is literally \textbf{a head}, exactly as in English \emph{head of department}.

The family is large: \textbf{el capítulo} (a chapter, a little head), \textbf{la capital} (the head city), \textbf{acabar} — to finish, to bring to a head — and English \textbf{captain}, \textbf{chief}, \textbf{chef}, \textbf{capital} and \textbf{cape}. A \emph{chef} is the head of a kitchen, and \emph{chief} and \emph{jefe} are the same word in two languages.
\end{cousinweb}

\subsection*{Your turn}
Say these aloud:

\begin{itemize}
\raggedright
  \item I have to speak to the boss
  \item the head of HR, using the feminine
  \item that you cannot make your mind up about asking her
  \item what \emph{jefe} and \emph{cabeza} have in common
\end{itemize}

\begin{tcolorbox}[breakable,colback=teal!4,colframe=teal!35!black,title={Before you move on}]
The boss? (\textbf{El jefe}.) The feminine? (\textbf{La jefa}.) What does \emph{caput} mean? (\textbf{Head} — the same one inside \emph{la cabeza}.)
\end{tcolorbox}
\section[(reducir, decidirse, el jefe)]{(reducir, decidirse, el jefe) — from cold}

\label{lesson:ES-R462-repaso-jefe}



\begin{quote}
Close the lessons before this one. Say all three: \emph{to cut back}, \emph{to make your mind up}, \emph{the boss}.
\end{quote}

\begin{grammarlens}[title={Grammar Lens: one word, two roads, four hundred years apart}]
\noindent
\begin{tabularx}{\linewidth}{@{}>{\raggedright\arraybackslash}X>{\raggedright\arraybackslash}X@{}}
\toprule
\textbf{the road} & \textbf{the word} \\
\midrule
worn down through Spanish speech & \textbf{la cabeza} \\
borrowed from French, much later & \textbf{el jefe} \\
\bottomrule
\end{tabularx}

Both are Latin \emph{caput}, head. The first went through centuries of mouths and came out as a body part; the second was picked up whole from a neighbour and came out as a person.

You have watched this several times — \emph{cosa} against \emph{causa}, \emph{razón} against \emph{ración}, \emph{rueda} against \emph{rotativo}, \emph{consejo} against \emph{consultar}. Every one of those came out of Latin twice at two different speeds. This one took a detour through a \textbf{third} language, which is why it looks like neither.
\end{grammarlens}

\begin{grammarlens}[title={Grammar Lens: two verbs, two pictures of cutting and leading}]
\noindent
\begin{tabularx}{\linewidth}{@{}>{\raggedright\arraybackslash}X>{\raggedright\arraybackslash}X@{}}
\toprule
\textbf{verb} & \textbf{the picture inside it} \\
\midrule
\textbf{reducir} & \emph{re-ducere}, to \textbf{lead back} \\
\textbf{decidirse} & \emph{de-caedere}, to \textbf{cut off} \\
\bottomrule
\end{tabularx}

Deciding cuts the other options away; reducing leads the number back down. Both are physical, and both are still doing the work: \emph{la ducha} leads water, \emph{preciso} is cut to the edge.

And \emph{decidirse} is the fourth \emph{-se} pair you have met, after \emph{mudarse}, \emph{acercarse} and \emph{arreglarse}. Here the \emph{-se} adds effort — \emph{no me decido} is dithering, where \emph{no decido} would be a job description.
\end{grammarlens}

\subsection*{Your turn}
Say these aloud:

\begin{itemize}
\raggedright
  \item the boss wants to cut costs
  \item I can't make my mind up
  \item which two Spanish words are both Latin \emph{caput}
\end{itemize}

\begin{tcolorbox}[breakable,colback=teal!4,colframe=teal!35!black,title={Before you move on}]
All three again, in Spanish, without looking. Then one sentence about a decision somebody is putting off.
\end{tcolorbox}
\section[Practice]{(el correo del jefe) — read it, then do it yourself}

\label{lesson:ES-C462-sintesis-jefe}



\begin{quote}
Say \textbf{el jefe} and \textbf{reducir}. The message below is from one about the other.
\end{quote}

\subsection*{Your turn}
An internal message, and one reply. Read both, then answer.

\textbf{1 — de la jefa}

\begin{quote}
Buenos días. Tenemos que \textbf{reducir} el gasto de material este trimestre.
\end{quote}

\begin{quote}
Antes del viernes, decidme qué se puede dejar de comprar en vuestra sección.
\end{quote}

\textbf{2 — la respuesta}

\begin{quote}
Lo miro hoy. Todavía no \textbf{me decido} entre quitar el papel de color o bajar los pedidos de tinta, así que te digo mañana.
\end{quote}

Say these aloud:

\begin{itemize}
\raggedright
  \item what has to be cut, and by when
  \item the two things the second person is choosing between
  \item when the answer is coming
\end{itemize}

\begin{culture}
\textbf{No me decido entre} is the phrase to keep whole. With \textbf{entre} you name the options you are stuck between; with \textbf{por} you name the one you finally take. \emph{No me decido entre A y B}, then \emph{me decidí por A}.

\textbf{Decidme} is \emph{vosotros}, plural and informal — a boss writing to her own team in Spain. A Latin American writer would say \emph{díganme}, and a formal Spanish notice to the public would too. The verb ending places the writer, the reader and the register in one word.

And \textbf{tenemos que reducir}: \emph{nosotros}, although the decision is hers. Spanish management writing does this constantly, and it is doing the same softening as the \emph{se} on a shop notice — spreading the action out so nobody is singled out as its author.
\end{culture}

\begin{tcolorbox}[breakable,colback=teal!4,colframe=teal!35!black,title={Before you move on}]
Now your turn, out loud and then in writing. Four lines: as a boss, say spending has to be cut and ask for suggestions by Friday; then as the reply, say you are torn between two options and will answer tomorrow.

Then check yourself: did you use \emph{entre} for the options and \emph{por} for the choice?
\end{tcolorbox}
